% ======================================================================
%  main.tex
%  Компактная шпаргалка: 14 x 16 см, 8 pt, pdfLaTeX.
%
%  Вся вёрстка, пакеты и макросы лежат в cheatsheet.sty.
%  Здесь остаются только метаданные и содержание.
%
%  Быстрый старт:
%    1. Замените \title, \author и вступительный текст.
%    2. Каждый вопрос оформляйте через \tquastion{...}.
%    3. Ответ пишите сразу после \tquastion.
%    4. Соберите: latexmk -pdf main.tex
%    5. Для печати 3 x 4 на A4: latexmk -pdf collage.tex
% ======================================================================
\documentclass[8pt]{extarticle}
\usepackage{cheatsheet}

% Раскомментируйте пакет, если нужен водяной знак.
% \usepackage{draftwatermark}

% ----------------------------------------------------------------------
%  Метаданные. Замените под свою дисциплину.
% ----------------------------------------------------------------------
\title{Шпаргалка по дисциплине \\ <<Название дисциплины>>}
\author{Имя Фамилия, группа}
\date{ }

% Водяной знак. Сначала раскомментируйте \usepackage{draftwatermark} выше.
% \SetWatermarkText{\textbf{Ваше имя}}
% \SetWatermarkColor[gray]{0.95}
% \SetWatermarkScale{1}

\begin{document}
	\maketitle

	% Вступительный текст. Замените его своим описанием.
	Краткое описание шпаргалки: дисциплина, преподаватель, семестр. Этот абзац нужен только для примера, его можно удалить.

	\tableofcontents

	% ==================================================================
	%  ПРИМЕР ВОПРОСА. Удалите его целиком и добавьте свои вопросы.
	% ==================================================================
	\tquastion{1. Пример вопроса. Замените этот текст.}

	% Ключевая формула в рамке. Внутри gathered можно делать переносы.
	\[\boxed{\begin{gathered}
		\text{\scriptsize\textbf{Ключевая формула}}\\[3pt]
		a^2+b^2=c^2
	\end{gathered}}\]

	% Тонкая линия и расшифровка обозначений, как в готовых ответах.
	\noindent\rule{\linewidth}{0.4pt}\par\smallskip
	\noindent\textit{Обозначения:} $a,b$ обозначают катеты, $c$ обозначает гипотенузу.\par\smallskip

	Этот абзац показывает, как выглядит основной текст ответа. Пишите коротко: определение, формула, пояснение и, если нужно, пример.

	\begin{itemize}
		\item Первый тезис ответа.
		\item Второй тезис ответа.
	\end{itemize}

	% Пример таблицы. Для широких таблиц используйте L{...} или C{...}.
	\begin{center}
		\begin{tabular}{@{}ll@{}}
			\toprule
			Величина & Значение \\
			\midrule
			$a$ & 3 \\
			$b$ & 4 \\
			$c$ & 5 \\
			\bottomrule
		\end{tabular}
	\end{center}

	% ==================================================================
	%  Шаблон следующего вопроса. Раскомментируйте и заполните.
	% ==================================================================
	% \tquastion{2. Название следующего вопроса.}
	%
	% \[\boxed{\begin{gathered}
	% 	\text{\scriptsize\textbf{Ключевая формула}}\\[3pt]
	% 	...
	% \end{gathered}}\]
	%
	% \noindent\rule{\linewidth}{0.4pt}\par\smallskip
	% \noindent\textit{Обозначения:} ...\par\smallskip
	%
	% Основной текст ответа.

	% ==================================================================
	%  Список источников. Раскомментируйте, если он нужен.
	% ==================================================================
	% \newpage
	% \addcontentsline{toc}{subsection}{Источники}
	% \begin{thebibliography}{9}
	% 	\bibitem{source1} Автор, название, год.
	% 	\bibitem{source2} Автор, название, год.
	% \end{thebibliography}

\end{document}
